\section{Related Work}
\label{sec:related}

Three literatures bear on the question of when organization is load-bearing.
Each has established a piece of it; the gap is stated at the end.

\paragraph{Role specialization and role drift.}
That prompted roles are not reliably obeyed is established. Cemri et
al.~\cite{cemri2025why} annotate 1{,}600+ traces across seven frameworks and
derive MAST, a 14-mode failure taxonomy in which disobeying the role
specification is a named mode. Cao et al.~\cite{cao2026modules} give the
training-time account: \emph{role drift}, in which modules ``preserve or
improve end-task performance while deviating from their assigned roles'' ---
86\% of one pipeline's apparent RL gain vanishes once the decomposer is held
to its role. Orogat et al.~\cite{orogat2026architectural} hold the LLM fixed
and vary architecture, and find that prompt-level role conditioning does not
induce the intended specialization. Patel~\cite{patel2026representational}
measures the phenomenon we call collapse directly, reporting near-identical
rationales across role-prompted committee members, and
Kandoussi~\cite{kandoussi2026whoami} reports the converse --- behavioral
differentiation arising \emph{without} assigned roles, and collapsing toward
uniformity when scaffolding is stripped. Enforcement-side work makes role
scope structural rather than advisory: Stettler et
al.~\cite{stettler2026progressive} deliver role-scoped tools through a server
that enforces scope deterministically, on the premise that prompt-defined
policy is probabilistic advice; Ge and Guo~\cite{ge2026roles} gate which edits
to a role pool may be committed, though the behavioral payload of an evolved
role remains a system prompt.

\emph{What is already owned here.} ``A role-quality metric can decouple from
task performance'' is not available as a contribution. Cao et
al.~\cite{cao2026modules} own it by name for training-time drift, and Yang et
al.~\cite{yang2026precise} own it for this architecture and one of these
actors: reviewer precision reaches $0.861$ while critique uptake does not
follow, on the same corpus this work instruments. We state that overlap
plainly rather than work around it, and we do not claim the decoupling itself.

\paragraph{Redundancy and pruning in multi-agent systems.}
A second literature asks whether the agents are needed at all. Chen et
al.~\cite{chen2026agentslimming} prune or downgrade redundant agents in
graph-structured workflows under a baseline-anchored acceptance rule, cutting
token cost by up to 78.9\%. Dylan et al.~\cite{dylan2026equal} hold total LLM
calls fixed and find an evolved Planner--Executor--Critic team no better than
an evolved single agent ($0.769$ vs.\ $0.754$, $p=0.80$), with a leave-one-in
analysis attributing the entire benefit to the executor. Yu et
al.~\cite{yu2026bottleneck} give an information-bottleneck account in which
multi-agent structure pays only under bounded relay bandwidth, and predict
that gains shrink or reverse for stronger models. Zhang et
al.~\cite{zhang2026rethinking} sharpen this into a critique of the
evaluations themselves: reported multi-agent gains, they argue, are largely
setup artifacts arising from \emph{structural collapse}, disabled tool paths,
or baselines that never demanded multi-agent structure, in settings where
random pruning already preserves accuracy.

\emph{What is already owned here, and the objection we must answer.}
``Agents are redundant'' is owned by this literature. More pointedly,
\cite{zhang2026rethinking} is aimed at exactly the observation we report for
gpt-oss --- collapse with no accuracy cost --- and would read it as evidence
that the setup never required multi-agent structure. This is the strongest
objection to our headline and we treat it as a test rather than a discussion
point: \S\ref{subsec:notsingle} reports a budget-matched single-agent control
on the same problems, and the collapsed team beats it. Collapse in our setting
is not reducible to a single agent, which is what their account would predict.

\paragraph{Coordination versus capability.}
The third literature asks who benefits from structure, and it is remarkably
consistent in direction. Canedo~\cite{canedo2026architecture} gives one coding
agent five informationally equivalent specification formats across six models
and reports a strong format~$\times$~model interaction: format barely matters
for the strongest models (spread $0.17$--$0.92$) and matters greatly for the
weakest ($0.83$--$2.42$), framing structured specification as a
\emph{capability equalizer} whose value is inversely proportional to model
strength. Yu et al.~\cite{yu2026bottleneck} reach the same monotone
prediction from theory. Dochkina~\cite{dochkina2026drop} argues the
complementary case and is the paper most directly opposed to our framing:
given minimal scaffolding, agents self-organize and outperform designed
structures, with autonomy's payoff \emph{increasing} in model capability.
Read together with the equalizer account, the field's prediction is a single
monotone ordering --- imposed structure helps weak models, autonomy suits
strong ones.

\emph{We address \cite{dochkina2026drop} head-on because half our result
resembles it and half contradicts it.} Our stronger actor does discard its
assigned organization and lose nothing, which is what a self-organization
account predicts. But it does not self-organize into anything: it collapses
to three byte-identical outputs (first-pass rate $0.698$), which is the
opposite of spontaneous differentiation, and its advantage over a single agent
survives that collapse rather than arising from any new structure. Meanwhile
the weaker actor, which the equalizer account says should depend on structure
most visibly, keeps its structure perfectly intact ($0.000$ first-pass
collapse) and loses accuracy anyway. Neither account predicts an actor pair
in which the organizational measure and the performance cost point in opposite
directions.

\paragraph{The gap.}
Prior work establishes that roles are often not obeyed, that agents are often
redundant, and that structure tends to help weaker models. What is missing is
a test of whether \emph{organization} is load-bearing that holds compute and
information fixed, measures the organization directly rather than inferring it
from outcome, intervenes in both directions --- removing role instructions and
enforcing them architecturally --- and runs on more than one actor substrate.
Doing so yields a result no monotone capability account predicts: the
\emph{more} capable actor loses its organization almost entirely and pays
nothing, the \emph{less} capable one retains it fully and pays anyway, and
restoring the organization by construction does not restore the performance.
The sign inversion, not the decoupling, is what we claim.

\paragraph{Studies of the same trace corpus.}
Three prior studies analyze the trajectory corpus this work instruments.
Jiang et al.~\cite{jiang2026confidence} audit a confidence-routed broadcast
protocol, separating routing, calibration, and commitment. Yang et
al.~\cite{yang2026protocol} study cost-aware routing \emph{across} protocols,
finding a probe predicts baseline failure well ($\mathrm{AUROC}=0.885$) but
predicts protocol-specific value poorly. Yang et al.~\cite{yang2026precise}
study reviewer precision and critique uptake within this architecture on one
of these actors. All three are observational, none intervenes on the
organization, and none measures whether the specified division of labor is
realized. The controlled arms of \S\ref{sec:method} do not exist in that
corpus, and they are what separates this work from a fourth reanalysis of it.
